\FloatBarrier
\section{A re-implementation of Latent-LoRA}
\label{app:latent}

Latent-LoRA \citep{latentlora2026} is the published method closest to our recognition step. It has no
public code that we could find. We wrote it again from the text of its paper and ran it next to the
certified switch. This was not planned in advance: no rule for this run was written in the plan
before it was started.

\paragraph{What was run.} The five tasks of Section~\ref{sec:five}, \LatentTrainInputs{} training and
\LatentTestInputs{} test inputs per task, T5-base, \LatentEpochs{} epochs, \LatentSeeds{} seeds (42 to 44), on two T4 GPUs.
Both methods ran in the same stand-alone code with the same data, budget and test. The certified switch
is the stand-alone code of \texttt{certified\_switch/} (rank \LatentRank{}); all-on is the same set of
adapters with every adapter used. For Latent-LoRA we ran the rule of the paper (``soft'': the
adapters are blended with the weights given by a router) and the ablation of the paper without router
(``sum''). The differences from the paper that we know of: \LatentEpochs{} epochs instead of 30, T5-base
instead of T5-large, five tasks, and our own Rouge-L scorer. The paper does not state 20 further details;
the choice made for each is listed in \texttt{latent\_lora/README.md} in the repository.

\begin{table}[h]
  \caption{Original instructions, mean over \LatentSeeds{} seeds. Answers changed: answers to old inputs that
  differ from those given right after learning. The sum rule and all-on use no router. Rows 3 and 4 are
  our re-implementation from the paper.}
  \label{tab:latentmain}
  \centering
  \small
  \begin{tabular}{lccccc}
    \toprule
    Method & \shortstack[c]{Score\\right after\\learning $\uparrow$} & \shortstack[c]{Final\\score $\uparrow$} & \shortstack[c]{Forgetting\\(points) $\downarrow$} & \shortstack[c]{Answers\\changed (of\\1{,}200) $\downarrow$} & \shortstack[c]{Sent to own\\task, \% $\uparrow$} \\
    \midrule
    \cert{} (ours) & 76.49 & 76.49 & 0.00 & 0 & 100.0 \\
    \plain{} (same adapters, all on) & 76.49 & 36.04 & 50.56 & 970 & -- \\
    \midrule
    Latent-LoRA, soft rule & 42.55 & 42.55 & 0.00 & 0 & 100.0 \\
    Latent-LoRA, sum rule & 48.71 & 47.66 & 1.31 & 429 & -- \\
    \bottomrule
\end{tabular}

\end{table}

\begin{table}[h]
  \caption{The four wordings of the task instructions (Appendix~\ref{app:rewordings}), mean over \LatentSeeds{} seeds. Answers changed: of \LatentCompared{} answers. Right: share of test inputs given to their own task (for the soft rule: the adapter with the largest router weight).}
  \label{tab:latentwordings}
  \centering
  \small
  \begin{tabular}{lcccccc}
    \toprule
    & \multicolumn{3}{c}{\cert{} (ours)} & \multicolumn{3}{c}{Latent-LoRA, soft rule} \\
    \cmidrule(lr){2-4} \cmidrule(lr){5-7}
    Wording & \shortstack[c]{Final\\score $\uparrow$} & \shortstack[c]{Answers\\changed $\downarrow$} & \shortstack[c]{Sent to own\\task, \% $\uparrow$} & \shortstack[c]{Final\\score $\uparrow$} & \shortstack[c]{Answers\\changed $\downarrow$} & \shortstack[c]{Sent to own\\task, \% $\uparrow$} \\
    \midrule
    original & 76.49 & 0 & 100.0 & 42.55 & 0 & 100.0 \\
    light & 75.87 & 0 & 100.0 & 48.71 & 0 & 100.0 \\
    medium & 75.91 & 0 & 100.0 & 32.53 & 0 & 100.0 \\
    heavy & 71.74 & 0 & 94.0 & 27.75 & 0 & 100.0 \\
    \bottomrule
\end{tabular}

\end{table}

\begin{table}[h]
  \caption{Original instructions, mean over \LatentSeeds{} seeds: score of each task right after it
  was learned, and the training loss of its last epoch.}
  \label{tab:latentpertask}
  \centering
  \small
  \begin{tabular}{lcccc}
    \toprule
    & \multicolumn{2}{c}{Score right after learning $\uparrow$} & \multicolumn{2}{c}{Training loss, last epoch $\downarrow$} \\
    \cmidrule(lr){2-3} \cmidrule(lr){4-5}
    Task (in order) & \cert{} (ours) & Latent-LoRA & \cert{} (ours) & Latent-LoRA \\
    \midrule
    relation extr. & 99.3 & 91.1 & 0.02 & 0.40 \\
    sst2 & 77.7 & 55.3 & 0.19 & 0.45 \\
    glucose & 59.5 & 2.0 & 0.61 & 2.02 \\
    sentiment140 & 72.0 & 37.0 & 0.23 & 0.74 \\
    emotion & 74.0 & 27.3 & 0.32 & 1.91 \\
    \bottomrule
\end{tabular}

\end{table}

\paragraph{Learning.} Our re-implementation learned the first task and almost nothing after it
(Table~\ref{tab:latentpertask}). With seed 42, the scores right after learning were
\LatentSeedSoftRelation{}, \LatentSeedSoftSst{}, \LatentSeedSoftGlucose{}, \LatentSeedSoftSenti{} and
\LatentSeedSoftEmotion{} for relation extraction, SST-2, Glucose, Sentiment140 and emotion; the other
seeds are alike. At the end of the Glucose task the training loss was about \LatentLossSoftGlucose{}
against about \LatentLossCertGlucose{} for the certified switch. Learned and AP are \LatentSoftAP{}
against \LatentCertAP{}. We do not know whether this comes from the reduced budget (a third of the
epochs, a smaller model) or from a mistake in our re-implementation, and without public code we cannot
check. So the run says nothing about how well Latent-LoRA learns or retains, and we make no such claim.
One sign that its adapters were barely trained: with the light rewording its average is higher
(\LatentSoftAPLight{}) than with the original instructions (\LatentSoftAP{}), because its Sentiment140
score rises from \LatentSoftSentOriginal{} with the original instruction to between \LatentSoftSentLightLow{} and \LatentSoftSentLightHigh{} (over seeds) with the
light rewording, and falls to \LatentSoftSentMediumMean{} with the medium rewording.

\paragraph{Its router.} The router of the re-implementation gave every test input to its own task in
all four wordings (right task \LatentSoftRight{}\%; certified switch \LatentCertRightOriginal{}\% /
\LatentCertRightLight{}\% / \LatentCertRightMedium{}\% / \LatentCertRightHeavy{}\% in the four wordings), and with
the soft rule none of \LatentCompared{} answers changed in any wording. So on these five tasks its
router is as stable as ours in practice, and it recognises the heavy rewording better than ours does.
What the certified switch adds is the proof and the
certificate on each answer, not better recognition. Latent-LoRA was not run in the stress test with a
look-alike task (Section~\ref{sec:stress}), so we do not know how its router behaves there. With the
sum rule (no router) AP is \LatentSumAP{}, FT \LatentSumFT{} and \LatentSumChanged{} of \LatentCompared{} answers changed.

\paragraph{The stand-alone code.} The same run reproduces the result of the certified switch on T5-base in a
second, independent code base: AP \LatentCertAP{} (\LatentSeeds{} seeds), FT \LatentCertFT{}, and
\LatentCertChanged{} of \LatentCompared{} answers changed in each of the four wordings (AP
\LatentCertAPLight{} / \LatentCertAPMedium{} / \LatentCertAPHeavy{} with the light, medium and heavy
rewording; right task \LatentCertRightHeavy{}\% with the heavy one). The same adapters all on give AP
\LatentAllAP{}, FT \LatentAllFT{} and \LatentAllChanged{} of \LatentCompared{} answers changed. The main
experiments of the paper use other code and another rank; their AP is \FiveCertAP{}.

\paragraph{What this does not show.} Nothing about the learning or the forgetting of Latent-LoRA as
its authors built it. Nothing about its router outside these five tasks. The reworded instructions
are the development set described in Appendix~\ref{app:rewordings}.
